% Source: https://artofproblemsolving.com/community/c3h2131005_volume_of_tetrahedron
% Stable identity: aops-topic-2131005-q01
Use the three edge vectors from $A$ (the sketch is schematic, not to scale):

\begin{asy}
usepackage("amsmath");
size(5cm);
defaultpen(fontsize(10pt));
pair A=(0,0), B=(1.6,2.8), C=(3.0,-0.6), D=(3.7,0.9);
draw(A--B--C--cycle,linewidth(0.9));
draw(B--D--C,linewidth(0.9));
draw(A--B,blue+linewidth(1.1),Arrow(7)); draw(A--C,blue+linewidth(1.1),Arrow(7)); draw(A--D,blue+dashed+linewidth(1.1),Arrow(7));
dot(A); dot(B); dot(C); dot(D);
label("$A(2,3,5)$",A,W,fontsize(8pt)); label("$B(-1,6,9)$",B,N,fontsize(8pt));
label("$C(-4,-10,-4)$",C,SE,fontsize(8pt)); label("$D(11,2,7)$",D,E,fontsize(8pt));
label("\mbox{}",(min(currentpicture,user=true)+max(currentpicture,user=true))/2,black); // keeps the text after the figure black when LaTeX typesets the labels inline
\end{asy}

\[\overrightarrow{AB}=(-3,3,4),\qquad\overrightarrow{AC}=(-6,-13,-9),\qquad\overrightarrow{AD}=(9,-1,2).\]
The parallelepiped spanned by three vectors has volume equal to the absolute value of the determinant of the matrix with those vectors as rows. A tetrahedron is a pyramid on one of the triangular faces, whose area is half a face of the parallelepiped, with the same height, so its volume is $\dfrac{1}{3}\cdot\dfrac{1}{2}=\dfrac{1}{6}$ of the parallelepiped's. Expanding along the first row,
\[\det\begin{pmatrix}-3&3&4\\-6&-13&-9\\9&-1&2\end{pmatrix}=-3\bigl((-13)(2)-(-9)(-1)\bigr)-3\bigl((-6)(2)-(-9)(9)\bigr)+4\bigl((-6)(-1)-(-13)(9)\bigr)\]
\[=-3(-35)-3(69)+4(123)=105-207+492=390 .\]
So the volume is
\[\dfrac{|390|}{6}=\boxed{65}.\]
